% Generated draft: requires mathematical and source review.
\documentclass[11pt]{article}
\usepackage[margin=25mm]{geometry}
\usepackage{fontspec}
\setmainfont{Latin Modern Roman}
\usepackage{amsmath,amssymb,mathtools,unicode-math}
\setmathfont{Latin Modern Math}
\usepackage{enumitem,needspace,xurl,hyperref}
\setlength{\parindent}{0pt}
\setlength{\parskip}{5pt}
\setlength{\emergencystretch}{3em}
\newcommand{\N}{\mathbb{N}}
\newcommand{\Z}{\mathbb{Z}}
\newcommand{\Q}{\mathbb{Q}}
\newcommand{\R}{\mathbb{R}}
\newcommand{\C}{\mathbb{C}}
\newcommand{\F}{\mathbb{F}}
\newcommand{\A}{\mathbb{A}}
\newcommand{\PP}{\mathbb{P}}
\newcommand{\E}{\mathbb{E}}
\newcommand{\eps}{\varepsilon}
\newcommand{\dd}{\,\mathrm{d}}
\begin{document}
\section*{UnsolvedMath ID: 30004238}
\textbf{Source problem code:} catalogue entry
\subsection{Definitions and mathematical statement}
A finite abstract simplicial complex on a finite set \(V\) is a collection \(K\subseteq 2^V\) containing the empty face and every singleton and closed under taking subsets. Its vertex set is denoted by \(V(K)\), and its dimension is \(\max_{\sigma\in K}(|\sigma|-1)\). For a nonnegative integer \(d\), a combinatorial \(d\)-sphere means a finite simplicial complex piecewise-linearly homeomorphic to the boundary of a \((d+1)\)-simplex. A complex \(S\) is a subcomplex of \(S'\) when every face of \(S\) is a face of \(S'\).

The question is whether, for every nonnegative integer \(d\) and every combinatorial \(d\)-sphere \(S\) with \(n=|V(S)|\geq d+3\), there exists a combinatorial \((d+1)\)-sphere \(S'\) satisfying
\[
V(S')=V(S),\qquad S\subseteq S'.
\]
Thus no additional vertices are permitted, and all faces of the given sphere must be retained.

The source does not explicitly specify the range of \(d\) or define ``combinatorial sphere''; the formulation above uses nonnegative dimensions and the standard piecewise-linear convention. The supplied dated assessment labels the problem partially solved and reports the extension for flag, stacked, and join spheres, while describing the general question as unresolved. These are supplied literature qualifications, not independently verified claims; the record gives no precise hypotheses or theorem references for those special cases.
\subsection{Short English statement}
Must every combinatorial \(d\)-sphere with at least \(d+3\) vertices occur as a subcomplex of a combinatorial \((d+1)\)-sphere using exactly the same vertices? The supplied curation reports special cases but leaves the general question unresolved; this statement-only draft requires human mathematical review.
\subsection{Sources}
\begin{itemize}
\item Local UnsolvedMath registry, permanent ID \(30004238\), title: Extending Combinatorial Spheres on the Same Vertex Set; source problem code \texttt{catalogue entry}. The record belongs to the Oberwolfach Reports Open Problems collection, records proposed year 2019, and names no proposer. Its supplied status is \texttt{partially\_solved}. Its statement label is \texttt{corrected\_verified}, with the supplied explanation that mathematical content was preserved while notation, typography, and sentence structure were repaired; this is registry metadata, not independent verification here.
\item Supplied report citation: \emph{Geometric, Algebraic, and Topological Combinatorics} (2020), Oberwolfach Reports, pp. 2395--2472, DOI 10.4171/owr/2019/39. \url{https://doi.org/10.4171/owr/2019/39}. The background describes these reports as records of mathematical workshops and their open-problem discussions. The supplied literature-source label is ``Oberwolfach report.''
\item Supplied literature-source label: ``Datta.'' \url{https://arxiv.org/abs/2006.04470}. No title, full author details, or theorem locator is supplied.
\item Literature assessment checked at \texttt{2026-08-22}: supplied dated curation, not independently verified. It reports that the same-vertex-set extension is proved for flag, stacked, and join spheres, while the general combinatorial-sphere question remains unresolved. The background calls this dated literature triage derived from the supplied OWR curation and explicitly requests independent verification of time-sensitive status and cited claims.
\end{itemize}
\end{document}
